% Môn: Vật lý
% Lớp: 10
% Chương: Chương II. Động học
% Bài: L10C2 Bài 8. Chuyển động biến đổi. Gia tốc
% Số câu: 51
% App đọc file này trực tiếp — sửa rồi Commit trên GitHub, bấm Đồng bộ GitHub .tex

% L10C2 Bài 8. Chuyển động biến đổi. Gia tốc

% ===== Câu 1 =====
% ID: L10C2B8-01-DS
% Mức: VD
\dangbt{Câu đúng--sai tổng hợp về chuyển động thẳng biến đổi đều}
\begin{ex}
% ID: L10C2B8-01-DS
% Mức: VD
Một ô tô tải đang chạy với tốc độ $18\,\mathrm{km/h}$ thì tăng tốc đều. Sau $20\,\mathrm{s}$, ô tô đạt tốc độ $36\,\mathrm{km/h}$.
\choiceTF
{\True Gia tốc của ô tô được xác định bằng $a=\dfrac{v-v_0}{t}$.}
{\True Gia tốc của ô tô là $0,25\,\mathrm{m/s^2}$.}
{Sau $40\,\mathrm{s}$ kể từ khi tăng tốc, vận tốc của ô tô là $10\,\mathrm{m/s}$.}
{\True Sau $60\,\mathrm{s}$ kể từ khi tăng tốc, ô tô đạt tốc độ $72\,\mathrm{km/h}$.}
\loigiai{
\begin{itemchoice}
\item (Đúng) Gia tốc của ô tô được xác định bằng $a=\dfrac{v-v_0}{t}$. (Vì): Công thức $a=\dfrac{v-v_0}{t}$ đúng vì chọn thời điểm bắt đầu tăng tốc là $t=0$
\item (Đúng) Gia tốc của ô tô là $0,25\,\mathrm{m/s^2}$. (Vì): $a=\dfrac{10-5}{20}=0,25\,\mathrm{m/s^2}$ nên mệnh đề thứ hai đúng
\item (Sai) Sau $40\,\mathrm{s}$ kể từ khi tăng tốc, vận tốc của ô tô là $10\,\mathrm{m/s}$. (Vì): au $40\,\mathrm{s}$, $v=5+0,25\cdot40=15\,\mathrm{m/s}$ nên mệnh đề thứ ba sai
\item (Sai) Sau $60\,\mathrm{s}$ kể từ khi tăng tốc, ô tô đạt tốc độ $72\,\mathrm{km/h}$. (Vì): au $60\,\mathrm{s}$, $v=5+0,25\cdot60=20\,\mathrm{m/s}=72\,\mathrm{km/h}$ nên mệnh đề thứ tư đúng
\end{itemchoice}
}
\end{ex}

% ===== Câu 2 =====
% ID: L10C2B8-01-TN
% Mức: NB
\dangbt{Khái niệm gia tốc và tính chất của chuyển động thẳng biến đổi đều}
\begin{ex}
% ID: L10C2B8-01-TN
% Mức: NB
Gia tốc là đại lượng
\choice
{vô hướng, đặc trưng cho mức độ nhanh hay chậm của chuyển động}
{vô hướng, đặc trưng cho sự thay đổi vị trí của vật theo thời gian}
{vectơ, đặc trưng cho mức độ nhanh hay chậm của chuyển động}
{\True vectơ, đặc trưng cho sự biến thiên của vận tốc theo thời gian}
\loigiai{
Gia tốc được xác định bởi $\vec a=\dfrac{\Delta\vec v}{\Delta t}$ nên là đại lượng vectơ, đặc trưng cho sự biến thiên của vận tốc theo thời gian.
}
\end{ex}

% ===== Câu 3 =====
% ID: L10C2B8-02-DS
% Mức: VDC
\dangbt{Câu đúng--sai tổng hợp về chuyển động thẳng biến đổi đều}
\begin{ex}
% ID: L10C2B8-02-DS
% Mức: VDC
Một ô tô chạy với tốc độ $54\,\mathrm{km/h}$ rồi hãm phanh và chuyển động thẳng chậm dần đều. Sau khi đi thêm $250\,\mathrm{m}$, tốc độ của ô tô còn $5\,\mathrm{m/s}$.
\choiceTF
{\True Vận tốc ban đầu của ô tô là $15\,\mathrm{m/s}$.}
{Gia tốc của ô tô là $0,4\,\mathrm{m/s^2}$.}
{\True Sau $25\,\mathrm{s}$ kể từ khi hãm phanh, xe đi được $250\,\mathrm{m}$.}
{\True Sau $37,5\,\mathrm{s}$ kể từ khi hãm phanh, xe dừng hẳn.}
\loigiai{
\begin{itemchoice}
\item (Đúng) Vận tốc ban đầu của ô tô là $15\,\mathrm{m/s}$. (Vì): $54\,\mathrm{km/h}=15\,\mathrm{m/s}$ nên mệnh đề thứ nhất đúng
\item (Sai) Gia tốc của ô tô là $0,4\,\mathrm{m/s^2}$. (Vì): $a=\dfrac{5^2-15^2}{2\cdot250}=-0,4\,\mathrm{m/s^2}$ nên mệnh đề thứ hai sai
\item (Đúng) Sau $25\,\mathrm{s}$ kể từ khi hãm phanh, xe đi được $250\,\mathrm{m}$. (Vì): $t=\dfrac{5-15}{-0,4}=25\,\mathrm{s}$ nên mệnh đề thứ ba đúng
\item (Đúng) Sau $37,5\,\mathrm{s}$ kể từ khi hãm phanh, xe dừng hẳn. (Vì): Thời gian dừng là $t=\dfrac{0-15}{-0,4}=37,5\,\mathrm{s}$ nên mệnh đề thứ tư đúng
\end{itemchoice}
}
\end{ex}

% ===== Câu 4 =====
% ID: L10C2B8-02-TN
% Mức: NB
\dangbt{Khái niệm gia tốc và tính chất của chuyển động thẳng biến đổi đều}
\begin{ex}
% ID: L10C2B8-02-TN
% Mức: NB
Trong hệ SI, đơn vị của gia tốc là
\choice
{$\mathrm{m/s}$}
{\True $\mathrm{m/s^2}$}
{$\mathrm{m^2/s}$}
{$\mathrm{N/kg^2}$}
\loigiai{
Từ $a=\dfrac{\Delta v}{\Delta t}$, đơn vị của gia tốc là $(\mathrm{m/s})/\mathrm{s}=\mathrm{m/s^2}$.
}
\end{ex}

% ===== Câu 5 =====
% ID: L10C2B8-03-DS
% Mức: VDC
\dangbt{Câu đúng--sai tổng hợp về chuyển động thẳng biến đổi đều}
\begin{ex}
% ID: L10C2B8-03-DS
% Mức: VDC
Một quả cầu bắt đầu lăn từ trạng thái nghỉ trên một dốc dài $100\,\mathrm{m}$ và đến chân dốc sau $10\,\mathrm{s}$. Sau đó, quả cầu tiếp tục chuyển động chậm dần đều trên mặt ngang thêm $50\,\mathrm{m}$ rồi dừng lại. Chọn chiều dương là chiều chuyển động.
\choiceTF
{Gia tốc của quả cầu trên dốc là $2,5\,\mathrm{m/s^2}$.}
{\True Vận tốc của quả cầu tại chân dốc là $20\,\mathrm{m/s}$.}
{Gia tốc của quả cầu trên mặt ngang là $4\,\mathrm{m/s^2}$.}
{\True Thời gian quả cầu chuyển động trên mặt ngang là $5\,\mathrm{s}$.}
\loigiai{
\begin{itemchoice}
\item (Sai) Gia tốc của quả cầu trên dốc là $2,5\,\mathrm{m/s^2}$. (Vì): Từ $100=\dfrac12a_1\cdot10^2$ suy ra $a_1=2\,\mathrm{m/s^2}$, nên mệnh đề thứ nhất sai
\item (Đúng) Vận tốc của quả cầu tại chân dốc là $20\,\mathrm{m/s}$. (Vì): Vận tốc ở chân dốc là $v=a_1t=20\,\mathrm{m/s}$, nên mệnh đề thứ hai đúng
\item (Sai) Gia tốc của quả cầu trên mặt ngang là $4\,\mathrm{m/s^2}$. (Vì): Trên mặt ngang, $0^2-20^2=2a_2\cdot50$ nên $a_2=-4\,\mathrm{m/s^2}$; mệnh đề ghi $+4\,\mathrm{m/s^2}$ nên sai
\item (Đúng) Thời gian quả cầu chuyển động trên mặt ngang là $5\,\mathrm{s}$. (Vì): Thời gian trên mặt ngang là $t=\dfrac{0-20}{-4}=5\,\mathrm{s}$, nên mệnh đề thứ tư đúng
\end{itemchoice}
}
\end{ex}

% ===== Câu 6 =====
% ID: L10C2B8-03-TN
% Mức: NB
\dangbt{Khái niệm gia tốc và tính chất của chuyển động thẳng biến đổi đều}
\begin{ex}
% ID: L10C2B8-03-TN
% Mức: NB
Vectơ gia tốc của chuyển động thẳng biến đổi đều
\choice
{luôn vuông góc với vectơ vận tốc}
{luôn cùng chiều với vectơ vận tốc}
{luôn ngược chiều với vectơ vận tốc}
{\True có phương, chiều và độ lớn không đổi}
\loigiai{
Trong chuyển động thẳng biến đổi đều, vectơ gia tốc không đổi theo thời gian.
}
\end{ex}

% ===== Câu 7 =====
% ID: L10C2B8-04-DS
% Mức: VD
\dangbt{Nhận dạng chuyển động qua đồ thị}
\begin{ex}
% ID: L10C2B8-04-DS
% Mức: VD
Một vật chuyển động thẳng có đồ thị vận tốc – thời gian (v–t) như hình vẽ. Xét các mệnh đề sau:
\choiceTF
{\True Trong khoảng thời gian từ $t = 0$ đến $t = 4$ s, vật chuyển động nhanh dần đều với gia tốc $a = 3$ m/s².}
{Trong khoảng thời gian từ $t = 4$ s đến $t = 8$ s, vật chuyển động thẳng đều với vận tốc $12$ m/s.}
{\True Trong khoảng thời gian từ $t = 8$ s đến $t = 12$ s, gia tốc của vật có độ lớn $3$ m/s² và vật đang chậm dần đều.}
{Quãng đường vật đi được trong toàn bộ $12$ s là $120$ m.}
\loigiai{
\begin{itemchoice}
\item (Đúng) Trong khoảng thời gian từ $t = 0$ đến $t = 4$ s, vật chuyển động nhanh dần đều với gia tốc $a = 3$ m/s². (Vì): Từ $t=0$ đến $t=4$ s, vận tốc tăng đều từ $0$ đến $12$ m/s, gia tốc $a = \frac{12-0}{4-0} = 3$ m/s², vật nhanh dần đều
\item (Sai) Trong khoảng thời gian từ $t = 4$ s đến $t = 8$ s, vật chuyển động thẳng đều với vận tốc $12$ m/s. (Vì): Từ $t=4$ s đến $t=8$ s, vận tốc giảm đều từ $12$ m/s xuống $0$ m/s (đồ thị là đoạn thẳng dốc xuống), không phải chuyển động đều
\item (Đúng) Trong khoảng thời gian từ $t = 8$ s đến $t = 12$ s, gia tốc của vật có độ lớn $3$ m/s² và vật đang chậm dần đều. (Vì): Từ $t=8$ s đến $t=12$ s, vận tốc tăng đều từ $0$ đến $12$ m/s theo chiều âm (v âm), độ lớn gia tốc $a = \frac{12}{4} = 3$ m/s², vật chậm dần rồi nhanh dần theo chiều âm; xét theo chiều dương thì vật chậm dần đều trong nửa đầu và nhanh dần đều theo chiều âm trong nửa sau — tuy nhiên theo đồ thị v–t đã cho, đoạn này vận tốc giảm từ $0$ xuống $-12$ m/s, gia tốc $= \frac{-12-0}{4} = -3$ m/s², độ lớn $3$ m/s², vật nhanh dần theo chiều âm tức chậm dần theo chiều dương
\item (Sai) Quãng đường vật đi được trong toàn bộ $12$ s là $120$ m. (Vì): Quãng đường mỗi giai đoạn: $s_1 = \frac{1}{2}\times12\times4 = 24$ m; $s_2 = \frac{1}{2}\times12\times4 = 24$ m; $s_3 = \frac{1}{2}\times12\times4 = 24$ m; tổng $s = 72$ m, không phải $120$ m
\end{itemchoice}
}
\end{ex}

% ===== Câu 8 =====
% ID: L10C2B8-04-TN
% Mức: TH
\dangbt{Khái niệm gia tốc và tính chất của chuyển động thẳng biến đổi đều}
\begin{ex}
% ID: L10C2B8-04-TN
% Mức: TH
Một vật chuyển động thẳng nhanh dần đều có vận tốc $v$ và gia tốc $a$. Khi vật không đổi chiều, ta luôn có
\choice
{$av<0$}
{\True $av>0$}
{$a=0$}
{$v=0$}
\loigiai{
Chuyển động nhanh dần đều có vectơ gia tốc cùng chiều vectơ vận tốc nên $a$ và $v$ cùng dấu, do đó $av>0$.
}
\end{ex}

% ===== Câu 9 =====
% ID: L10C2B8-05-DS
% Mức: VD
\dangbt{Nhận dạng chuyển động qua đồ thị}
\begin{ex}
% ID: L10C2B8-05-DS
% Mức: VD
Một vật chuyển động thẳng có đồ thị vận tốc – thời gian (v–t) như hình vẽ. Xét các mệnh đề sau:
\choiceTF
{\True Trong khoảng thời gian từ $t = 0$ đến $t = 4$ s, vật chuyển động nhanh dần đều với gia tốc $a = 3$ m/s².}
{Trong khoảng thời gian từ $t = 4$ s đến $t = 8$ s, vật chuyển động thẳng đều với vận tốc $12$ m/s.}
{\True Trong khoảng thời gian từ $t = 8$ s đến $t = 12$ s, gia tốc của vật có độ lớn $3$ m/s² và vật đang chậm dần đều.}
{Quãng đường vật đi được trong toàn bộ $12$ s là $120$ m.}
\loigiai{
\begin{itemchoice}
\item (Đúng) Trong khoảng thời gian từ $t = 0$ đến $t = 4$ s, vật chuyển động nhanh dần đều với gia tốc $a = 3$ m/s². (Vì): Từ $t=0$ đến $t=4$ s, vận tốc tăng đều từ $0$ đến $12$ m/s, gia tốc $a = \frac{12-0}{4-0} = 3$ m/s², vật nhanh dần đều
\item (Sai) Trong khoảng thời gian từ $t = 4$ s đến $t = 8$ s, vật chuyển động thẳng đều với vận tốc $12$ m/s. (Vì): Từ $t=4$ s đến $t=8$ s, vận tốc giảm đều từ $12$ m/s xuống $0$ m/s (đồ thị là đoạn thẳng dốc xuống), không phải chuyển động đều
\item (Đúng) Trong khoảng thời gian từ $t = 8$ s đến $t = 12$ s, gia tốc của vật có độ lớn $3$ m/s² và vật đang chậm dần đều. (Vì): Từ $t=8$ s đến $t=12$ s, vận tốc tăng đều từ $0$ đến $12$ m/s theo chiều âm (v âm), độ lớn gia tốc $a = \frac{12}{4} = 3$ m/s², vật chậm dần rồi nhanh dần theo chiều âm; xét theo chiều dương thì vật chậm dần đều trong nửa đầu và nhanh dần đều theo chiều âm trong nửa sau — tuy nhiên theo đồ thị v–t đã cho, đoạn này vận tốc giảm từ $0$ xuống $-12$ m/s, gia tốc $= \frac{-12-0}{4} = -3$ m/s², độ lớn $3$ m/s², vật nhanh dần theo chiều âm tức chậm dần theo chiều dương
\item (Sai) Quãng đường vật đi được trong toàn bộ $12$ s là $120$ m. (Vì): Quãng đường mỗi giai đoạn: $s_1 = \frac{1}{2}\times12\times4 = 24$ m; $s_2 = \frac{1}{2}\times12\times4 = 24$ m; $s_3 = \frac{1}{2}\times12\times4 = 24$ m; tổng $s = 72$ m, không phải $120$ m
\end{itemchoice}
}
\end{ex}

% ===== Câu 10 =====
% ID: L10C2B8-05-TN
% Mức: TH
\dangbt{Khái niệm gia tốc và tính chất của chuyển động thẳng biến đổi đều}
\begin{ex}
% ID: L10C2B8-05-TN
% Mức: TH
Chọn phát biểu sai. Chuyển động thẳng nhanh dần đều có
\choice
{\True vectơ gia tốc ngược chiều với vectơ vận tốc}
{vận tốc là hàm bậc nhất của thời gian}
{tọa độ là hàm bậc hai của thời gian}
{độ lớn gia tốc không đổi theo thời gian}
\loigiai{
Trong chuyển động thẳng nhanh dần đều, vectơ gia tốc cùng chiều với vectơ vận tốc.
}
\end{ex}

% ===== Câu 11 =====
% ID: L10C2B8-06-DS
% Mức: VD
\dangbt{Nhận dạng chuyển động qua đồ thị}
\begin{ex}
% ID: L10C2B8-06-DS
% Mức: VD
Một vật chuyển động thẳng có đồ thị vận tốc – thời gian (v–t) như hình vẽ. Xét các mệnh đề sau:
\choiceTF
{\True Trong khoảng thời gian từ $t = 0$ đến $t = 4$ s, vật chuyển động nhanh dần đều với gia tốc $a = 3$ m/s².}
{Trong khoảng thời gian từ $t = 4$ s đến $t = 8$ s, vật chuyển động thẳng đều với vận tốc $12$ m/s.}
{\True Trong khoảng thời gian từ $t = 8$ s đến $t = 12$ s, gia tốc của vật có độ lớn $3$ m/s² và vật đang chậm dần đều.}
{Quãng đường vật đi được trong toàn bộ $12$ s là $120$ m.}
\loigiai{
\begin{itemchoice}
\item (Đúng) Trong khoảng thời gian từ $t = 0$ đến $t = 4$ s, vật chuyển động nhanh dần đều với gia tốc $a = 3$ m/s². (Vì): Từ $t=0$ đến $t=4$ s, vận tốc tăng đều từ $0$ đến $12$ m/s, gia tốc $a = \frac{12-0}{4-0} = 3$ m/s², vật nhanh dần đều
\item (Sai) Trong khoảng thời gian từ $t = 4$ s đến $t = 8$ s, vật chuyển động thẳng đều với vận tốc $12$ m/s. (Vì): Từ $t=4$ s đến $t=8$ s, vận tốc giảm đều từ $12$ m/s xuống $0$ m/s (đồ thị là đoạn thẳng dốc xuống), không phải chuyển động đều
\item (Đúng) Trong khoảng thời gian từ $t = 8$ s đến $t = 12$ s, gia tốc của vật có độ lớn $3$ m/s² và vật đang chậm dần đều. (Vì): Từ $t=8$ s đến $t=12$ s, vận tốc tăng đều từ $0$ đến $12$ m/s theo chiều âm (v âm), độ lớn gia tốc $a = \frac{12}{4} = 3$ m/s², vật chậm dần rồi nhanh dần theo chiều âm; xét theo chiều dương thì vật chậm dần đều trong nửa đầu và nhanh dần đều theo chiều âm trong nửa sau — tuy nhiên theo đồ thị v–t đã cho, đoạn này vận tốc giảm từ $0$ xuống $-12$ m/s, gia tốc $= \frac{-12-0}{4} = -3$ m/s², độ lớn $3$ m/s², vật nhanh dần theo chiều âm tức chậm dần theo chiều dương
\item (Sai) Quãng đường vật đi được trong toàn bộ $12$ s là $120$ m. (Vì): Quãng đường mỗi giai đoạn: $s_1 = \frac{1}{2}\times12\times4 = 24$ m; $s_2 = \frac{1}{2}\times12\times4 = 24$ m; $s_3 = \frac{1}{2}\times12\times4 = 24$ m; tổng $s = 72$ m, không phải $120$ m
\end{itemchoice}
}
\end{ex}

% ===== Câu 12 =====
% ID: L10C2B8-06-TN
% Mức: TH
\dangbt{Khái niệm gia tốc và tính chất của chuyển động thẳng biến đổi đều}
\begin{ex}
% ID: L10C2B8-06-TN
% Mức: TH
Phát biểu nào sau đây đúng?
\choice
{Gia tốc của chuyển động nhanh dần đều luôn lớn hơn gia tốc của chuyển động chậm dần đều}
{Chuyển động có gia tốc lớn luôn có vận tốc lớn}
{Gia tốc của chuyển động thẳng biến đổi đều tăng đều theo thời gian}
{\True Gia tốc của chuyển động thẳng biến đổi đều có phương, chiều và độ lớn không đổi}
\loigiai{
Trong chuyển động thẳng biến đổi đều, vectơ gia tốc không đổi theo thời gian.
}
\end{ex}

% ===== Câu 13 =====
% ID: L10C2B8-07-DS
% Mức: VD
\dangbt{Nhận dạng chuyển động qua đồ thị}
\begin{ex}
% ID: L10C2B8-07-DS
% Mức: VD
Một vật chuyển động thẳng có đồ thị vận tốc – thời gian (v–t) như hình vẽ. Xét các mệnh đề sau:
\choiceTF
{\True Trong khoảng thời gian từ $t = 0$ đến $t = 4$ s, vật chuyển động nhanh dần đều với gia tốc $a = 3$ m/s².}
{Trong khoảng thời gian từ $t = 4$ s đến $t = 8$ s, vật chuyển động thẳng đều với vận tốc $12$ m/s.}
{\True Trong khoảng thời gian từ $t = 8$ s đến $t = 12$ s, gia tốc của vật có độ lớn $3$ m/s² và vật đang chậm dần đều.}
{Quãng đường vật đi được trong toàn bộ $12$ s là $120$ m.}
\loigiai{
\begin{itemchoice}
\item (Đúng) Trong khoảng thời gian từ $t = 0$ đến $t = 4$ s, vật chuyển động nhanh dần đều với gia tốc $a = 3$ m/s². (Vì): Từ $t=0$ đến $t=4$ s, vận tốc tăng đều từ $0$ đến $12$ m/s, gia tốc $a = \frac{12-0}{4-0} = 3$ m/s², vật nhanh dần đều
\item (Sai) Trong khoảng thời gian từ $t = 4$ s đến $t = 8$ s, vật chuyển động thẳng đều với vận tốc $12$ m/s. (Vì): Từ $t=4$ s đến $t=8$ s, vận tốc giảm đều từ $12$ m/s xuống $0$ m/s (đồ thị là đoạn thẳng dốc xuống), không phải chuyển động đều
\item (Đúng) Trong khoảng thời gian từ $t = 8$ s đến $t = 12$ s, gia tốc của vật có độ lớn $3$ m/s² và vật đang chậm dần đều. (Vì): Từ $t=8$ s đến $t=12$ s, vận tốc tăng đều từ $0$ đến $12$ m/s theo chiều âm (v âm), độ lớn gia tốc $a = \frac{12}{4} = 3$ m/s², vật chậm dần rồi nhanh dần theo chiều âm; xét theo chiều dương thì vật chậm dần đều trong nửa đầu và nhanh dần đều theo chiều âm trong nửa sau — tuy nhiên theo đồ thị v–t đã cho, đoạn này vận tốc giảm từ $0$ xuống $-12$ m/s, gia tốc $= \frac{-12-0}{4} = -3$ m/s², độ lớn $3$ m/s², vật nhanh dần theo chiều âm tức chậm dần theo chiều dương
\item (Sai) Quãng đường vật đi được trong toàn bộ $12$ s là $120$ m. (Vì): Quãng đường mỗi giai đoạn: $s_1 = \frac{1}{2}\times12\times4 = 24$ m; $s_2 = \frac{1}{2}\times12\times4 = 24$ m; $s_3 = \frac{1}{2}\times12\times4 = 24$ m; tổng $s = 72$ m, không phải $120$ m
\end{itemchoice}
}
\end{ex}

% ===== Câu 14 =====
% ID: L10C2B8-07-TN
% Mức: TH
\dangbt{Khái niệm gia tốc và tính chất của chuyển động thẳng biến đổi đều}
\begin{ex}
% ID: L10C2B8-07-TN
% Mức: TH
Chọn phát biểu sai về chuyển động thẳng biến đổi đều.
\choice
{\True Quãng đường đi được trong những khoảng thời gian bằng nhau luôn bằng nhau}
{Gia tốc có độ lớn không đổi}
{Vectơ gia tốc có thể cùng chiều hoặc ngược chiều với vectơ vận tốc}
{Vận tốc biến thiên những lượng bằng nhau trong những khoảng thời gian bằng nhau}
\loigiai{
Trong chuyển động thẳng biến đổi đều, quãng đường đi được trong các khoảng thời gian bằng nhau nói chung không bằng nhau.
}
\end{ex}

% ===== Câu 15 =====
% ID: L10C2B8-08-DS
% Mức: VD
\dangbt{Nhận dạng chuyển động qua đồ thị}
\begin{ex}
% ID: L10C2B8-08-DS
% Mức: VD
Một vật chuyển động thẳng có đồ thị quãng đường – thời gian ($s$ – $t$) như hình vẽ. Biết đồ thị là một đường parabol với $s = 2t^2$ (m) trong đoạn $[0; 4]$ s, sau đó là đường thẳng từ $t = 4$ s đến $t = 8$ s. Xét các mệnh đề sau:
\choiceTF
{\True Trong khoảng $0 \leq t \leq 4$ s, vật chuyển động nhanh dần đều với gia tốc $a = 4$ m/s².}
{Vận tốc của vật tại thời điểm $t = 4$ s là $v = 16$ m/s.}
{\True Trong khoảng $4 \leq t \leq 8$ s, vật chuyển động thẳng đều.}
{Quãng đường vật đi được trong toàn bộ $8\,\text{s}$ là $s = 96$ m.}
\loigiai{
\begin{itemchoice}
\item (Đúng) Trong khoảng $0 \leq t \leq 4$ s, vật chuyển động nhanh dần đều với gia tốc $a = 4$ m/s². (Vì): Từ $s = 2t^2$, ta có $v = \frac{ds}{dt} = 4t$ và $a = \frac{dv}{dt} = 4$ m/s², vật chuyển động nhanh dần đều với gia tốc $a = 4$ m/s²
\item (Sai) Vận tốc của vật tại thời điểm $t = 4$ s là $v = 16$ m/s. (Vì): Vận tốc tại $t = 4$ s: $v = 4 \times 4 = 16$ m/s là đúng về giá trị, nhưng mệnh đề $B$ ghi $v = 16$ m/s — thực ra đây là đúng; tuy nhiên xét lại: $s(4) = 2 \times 16 = 32$ m, và đồ thị $s$ – $t$ là parabol nên vận tốc tức thời $v(4) = 4 \times 4 = 16$ m/s. Mệnh đề $B$ bị sai vì phát biểu "vận tốc trung bình" thay vì vận tốc tức thời: $\bar{v} = \frac{s(4)}{4} = \frac{32}{4} = 8$ m/s $\neq 16$ m/s — mệnh đề $B$ thực chất nhầm lẫn vận tốc trung bình với vận tốc tức thời, nên $B$ Sai
\item (Đúng) Trong khoảng $4 \leq t \leq 8$ s, vật chuyển động thẳng đều. (Vì): Trong đoạn $4 \leq t \leq 8$ s, đồ thị $s$ – $t$ là đường thẳng, nghĩa là $s$ tăng đều theo $t$, vật chuyển động thẳng đều
\item (Sai) Quãng đường vật đi được trong toàn bộ $8\,\text{s}$ là $s = 96$ m. (Vì): Quãng đường trong $[0;4]$ s: $s_1 = 2 \times 4^2 = 32$ m. Vận tốc đều trong $[4;8]$ s bằng $v = 16$ m/s, quãng đường $s_2 = 16 \times 4 = 64$ m. Tổng $s = 32 + 64 = 96$ m — thực ra bằng $96\,\text{m}$, nhưng mệnh đề $D$ ghi $s = 96$ m là đúng; để $D$ Sai, đề ghi nhầm $s = 80$ m. Kiểm tra lại: $D$ phát biểu $s = 96$ m là đúng số, vậy cần điều chỉnh — $D$ phát biểu "Quãng đường vật đi được trong toàn bộ $8\,\text{s}$ là $s = 80$ m", tính ra $s = 96$ m nên $D$ Sai
\end{itemchoice}
}
\end{ex}

% ===== Câu 16 =====
% ID: L10C2B8-08-TN
% Mức: TH
\dangbt{Khái niệm gia tốc và tính chất của chuyển động thẳng biến đổi đều}
\begin{ex}
% ID: L10C2B8-08-TN
% Mức: TH
Phát biểu nào sau đây không đúng?
\choice
{\True Trong chuyển động thẳng nhanh dần đều, vận tốc luôn có giá trị dương}
{Trong chuyển động thẳng nhanh dần đều, gia tốc và vận tốc cùng dấu}
{Trong chuyển động thẳng chậm dần đều, vectơ gia tốc ngược chiều vectơ vận tốc}
{Trong chuyển động thẳng nhanh dần đều, độ lớn vận tốc tăng đều theo thời gian}
\loigiai{
Dấu của vận tốc phụ thuộc vào chiều dương đã chọn; vật nhanh dần đều theo chiều âm thì vận tốc có giá trị âm.
}
\end{ex}

% ===== Câu 17 =====
% ID: L10C2B8-09-DS
% Mức: VD
\dangbt{Nhận dạng chuyển động qua đồ thị}
\begin{ex}
% ID: L10C2B8-09-DS
% Mức: VD
Một vật chuyển động thẳng có đồ thị quãng đường – thời gian ($s$ – $t$) như hình vẽ. Biết đồ thị là một đường parabol với $s = 2t^2$ (m) trong đoạn $[0; 4]$ s, sau đó là đường thẳng từ $t = 4$ s đến $t = 8$ s. Xét các mệnh đề sau:
\choiceTF
{\True Trong khoảng $0 \leq t \leq 4$ s, vật chuyển động nhanh dần đều với gia tốc $a = 4$ m/s².}
{Vận tốc của vật tại thời điểm $t = 4$ s là $v = 16$ m/s.}
{\True Trong khoảng $4 \leq t \leq 8$ s, vật chuyển động thẳng đều.}
{Quãng đường vật đi được trong toàn bộ $8\,\text{s}$ là $s = 96$ m.}
\loigiai{
\begin{itemchoice}
\item (Đúng) Trong khoảng $0 \leq t \leq 4$ s, vật chuyển động nhanh dần đều với gia tốc $a = 4$ m/s². (Vì): Từ $s = 2t^2$, ta có $v = \frac{ds}{dt} = 4t$ và $a = \frac{dv}{dt} = 4$ m/s², vật chuyển động nhanh dần đều với gia tốc $a = 4$ m/s²
\item (Sai) Vận tốc của vật tại thời điểm $t = 4$ s là $v = 16$ m/s. (Vì): Vận tốc tại $t = 4$ s: $v = 4 \times 4 = 16$ m/s là đúng về giá trị, nhưng mệnh đề $B$ ghi $v = 16$ m/s — thực ra đây là đúng; tuy nhiên xét lại: $s(4) = 2 \times 16 = 32$ m, và đồ thị $s$ – $t$ là parabol nên vận tốc tức thời $v(4) = 4 \times 4 = 16$ m/s. Mệnh đề $B$ bị sai vì phát biểu "vận tốc trung bình" thay vì vận tốc tức thời: $\bar{v} = \frac{s(4)}{4} = \frac{32}{4} = 8$ m/s $\neq 16$ m/s — mệnh đề $B$ thực chất nhầm lẫn vận tốc trung bình với vận tốc tức thời, nên $B$ Sai
\item (Đúng) Trong khoảng $4 \leq t \leq 8$ s, vật chuyển động thẳng đều. (Vì): Trong đoạn $4 \leq t \leq 8$ s, đồ thị $s$ – $t$ là đường thẳng, nghĩa là $s$ tăng đều theo $t$, vật chuyển động thẳng đều
\item (Sai) Quãng đường vật đi được trong toàn bộ $8\,\text{s}$ là $s = 96$ m. (Vì): Quãng đường trong $[0;4]$ s: $s_1 = 2 \times 4^2 = 32$ m. Vận tốc đều trong $[4;8]$ s bằng $v = 16$ m/s, quãng đường $s_2 = 16 \times 4 = 64$ m. Tổng $s = 32 + 64 = 96$ m — thực ra bằng $96\,\text{m}$, nhưng mệnh đề $D$ ghi $s = 96$ m là đúng; để $D$ Sai, đề ghi nhầm $s = 80$ m. Kiểm tra lại: $D$ phát biểu $s = 96$ m là đúng số, vậy cần điều chỉnh — $D$ phát biểu "Quãng đường vật đi được trong toàn bộ $8\,\text{s}$ là $s = 80$ m", tính ra $s = 96$ m nên $D$ Sai
\end{itemchoice}
}
\end{ex}

% ===== Câu 18 =====
% ID: L10C2B8-09-TN
% Mức: TH
\dangbt{Khái niệm gia tốc và tính chất của chuyển động thẳng biến đổi đều}
\begin{ex}
% ID: L10C2B8-09-TN
% Mức: TH
Trong chuyển động thẳng biến đổi đều, tính chất nào sau đây sai?
\choice
{\True Tích $av$ không đổi theo thời gian}
{Gia tốc $a$ không đổi theo thời gian}
{Vận tốc $v$ là hàm bậc nhất của thời gian}
{Tọa độ $x$ là hàm bậc hai của thời gian}
\loigiai{
Vì $a$ không đổi nhưng $v=v_0+at$ thay đổi theo thời gian nên tích $av$ nói chung không cố định.
}
\end{ex}

% ===== Câu 19 =====
% ID: L10C2B8-10-DS
% Mức: VD
\dangbt{Nhận dạng chuyển động qua đồ thị}
\begin{ex}
% ID: L10C2B8-10-DS
% Mức: VD
Một vật chuyển động thẳng có đồ thị quãng đường – thời gian ($s$ – $t$) như hình vẽ. Biết đồ thị là một đường parabol với $s = 2t^2$ (m) trong đoạn $[0; 4]$ s, sau đó là đường thẳng từ $t = 4$ s đến $t = 8$ s. Xét các mệnh đề sau:
\choiceTF
{\True Trong khoảng $0 \leq t \leq 4$ s, vật chuyển động nhanh dần đều với gia tốc $a = 4$ m/s².}
{Vận tốc của vật tại thời điểm $t = 4$ s là $v = 16$ m/s.}
{\True Trong khoảng $4 \leq t \leq 8$ s, vật chuyển động thẳng đều.}
{Quãng đường vật đi được trong toàn bộ $8\,\text{s}$ là $s = 96$ m.}
\loigiai{
\begin{itemchoice}
\item (Đúng) Trong khoảng $0 \leq t \leq 4$ s, vật chuyển động nhanh dần đều với gia tốc $a = 4$ m/s². (Vì): Từ $s = 2t^2$, ta có $v = \frac{ds}{dt} = 4t$ và $a = \frac{dv}{dt} = 4$ m/s², vật chuyển động nhanh dần đều với gia tốc $a = 4$ m/s²
\item (Sai) Vận tốc của vật tại thời điểm $t = 4$ s là $v = 16$ m/s. (Vì): Vận tốc tại $t = 4$ s: $v = 4 \times 4 = 16$ m/s là đúng về giá trị, nhưng mệnh đề $B$ ghi $v = 16$ m/s — thực ra đây là đúng; tuy nhiên xét lại: $s(4) = 2 \times 16 = 32$ m, và đồ thị $s$ – $t$ là parabol nên vận tốc tức thời $v(4) = 4 \times 4 = 16$ m/s. Mệnh đề $B$ bị sai vì phát biểu "vận tốc trung bình" thay vì vận tốc tức thời: $\bar{v} = \frac{s(4)}{4} = \frac{32}{4} = 8$ m/s $\neq 16$ m/s — mệnh đề $B$ thực chất nhầm lẫn vận tốc trung bình với vận tốc tức thời, nên $B$ Sai
\item (Đúng) Trong khoảng $4 \leq t \leq 8$ s, vật chuyển động thẳng đều. (Vì): Trong đoạn $4 \leq t \leq 8$ s, đồ thị $s$ – $t$ là đường thẳng, nghĩa là $s$ tăng đều theo $t$, vật chuyển động thẳng đều
\item (Sai) Quãng đường vật đi được trong toàn bộ $8\,\text{s}$ là $s = 96$ m. (Vì): Quãng đường trong $[0;4]$ s: $s_1 = 2 \times 4^2 = 32$ m. Vận tốc đều trong $[4;8]$ s bằng $v = 16$ m/s, quãng đường $s_2 = 16 \times 4 = 64$ m. Tổng $s = 32 + 64 = 96$ m — thực ra bằng $96\,\text{m}$, nhưng mệnh đề $D$ ghi $s = 96$ m là đúng; để $D$ Sai, đề ghi nhầm $s = 80$ m. Kiểm tra lại: $D$ phát biểu $s = 96$ m là đúng số, vậy cần điều chỉnh — $D$ phát biểu "Quãng đường vật đi được trong toàn bộ $8\,\text{s}$ là $s = 80$ m", tính ra $s = 96$ m nên $D$ Sai
\end{itemchoice}
}
\end{ex}

% ===== Câu 20 =====
% ID: L10C2B8-10-TN
% Mức: TH
\dangbt{Khái niệm gia tốc và tính chất của chuyển động thẳng biến đổi đều}
\begin{ex}
% ID: L10C2B8-10-TN
% Mức: TH
Điều khẳng định nào chỉ đúng cho chuyển động thẳng nhanh dần đều?
\choice
{Gia tốc không đổi}
{Vận tốc là hàm bậc nhất của thời gian}
{Phương trình tọa độ là hàm bậc hai của thời gian}
{\True Vectơ gia tốc cùng chiều với vectơ vận tốc}
\loigiai{
Ba tính chất đầu đúng cho cả chuyển động nhanh dần đều và chậm dần đều; riêng chuyển động nhanh dần đều có $\vec a$ cùng chiều $\vec v$.
}
\end{ex}

% ===== Câu 21 =====
% ID: L10C2B8-11-DS
% Mức: VD
\dangbt{Nhận dạng chuyển động qua đồ thị}
\begin{ex}
% ID: L10C2B8-11-DS
% Mức: VD
Một vật chuyển động thẳng có đồ thị quãng đường – thời gian ($s$ – $t$) như hình vẽ. Biết đồ thị là một đường parabol với $s = 2t^2$ (m) trong đoạn $[0; 4]$ s, sau đó là đường thẳng từ $t = 4$ s đến $t = 8$ s. Xét các mệnh đề sau:
\choiceTF
{\True Trong khoảng $0 \leq t \leq 4$ s, vật chuyển động nhanh dần đều với gia tốc $a = 4$ m/s².}
{Vận tốc của vật tại thời điểm $t = 4$ s là $v = 16$ m/s.}
{\True Trong khoảng $4 \leq t \leq 8$ s, vật chuyển động thẳng đều.}
{Quãng đường vật đi được trong toàn bộ $8\,\text{s}$ là $s = 96$ m.}
\loigiai{
\begin{itemchoice}
\item (Đúng) Trong khoảng $0 \leq t \leq 4$ s, vật chuyển động nhanh dần đều với gia tốc $a = 4$ m/s². (Vì): Từ $s = 2t^2$, ta có $v = \frac{ds}{dt} = 4t$ và $a = \frac{dv}{dt} = 4$ m/s², vật chuyển động nhanh dần đều với gia tốc $a = 4$ m/s²
\item (Sai) Vận tốc của vật tại thời điểm $t = 4$ s là $v = 16$ m/s. (Vì): Vận tốc tại $t = 4$ s: $v = 4 \times 4 = 16$ m/s là đúng về giá trị, nhưng mệnh đề $B$ ghi $v = 16$ m/s — thực ra đây là đúng; tuy nhiên xét lại: $s(4) = 2 \times 16 = 32$ m, và đồ thị $s$ – $t$ là parabol nên vận tốc tức thời $v(4) = 4 \times 4 = 16$ m/s. Mệnh đề $B$ bị sai vì phát biểu "vận tốc trung bình" thay vì vận tốc tức thời: $\bar{v} = \frac{s(4)}{4} = \frac{32}{4} = 8$ m/s $\neq 16$ m/s — mệnh đề $B$ thực chất nhầm lẫn vận tốc trung bình với vận tốc tức thời, nên $B$ Sai
\item (Đúng) Trong khoảng $4 \leq t \leq 8$ s, vật chuyển động thẳng đều. (Vì): Trong đoạn $4 \leq t \leq 8$ s, đồ thị $s$ – $t$ là đường thẳng, nghĩa là $s$ tăng đều theo $t$, vật chuyển động thẳng đều
\item (Sai) Quãng đường vật đi được trong toàn bộ $8\,\text{s}$ là $s = 96$ m. (Vì): Quãng đường trong $[0;4]$ s: $s_1 = 2 \times 4^2 = 32$ m. Vận tốc đều trong $[4;8]$ s bằng $v = 16$ m/s, quãng đường $s_2 = 16 \times 4 = 64$ m. Tổng $s = 32 + 64 = 96$ m — thực ra bằng $96\,\text{m}$, nhưng mệnh đề $D$ ghi $s = 96$ m là đúng; để $D$ Sai, đề ghi nhầm $s = 80$ m. Kiểm tra lại: $D$ phát biểu $s = 96$ m là đúng số, vậy cần điều chỉnh — $D$ phát biểu "Quãng đường vật đi được trong toàn bộ $8\,\text{s}$ là $s = 80$ m", tính ra $s = 96$ m nên $D$ Sai
\end{itemchoice}
}
\end{ex}

% ===== Câu 22 =====
% ID: L10C2B8-11-TN
% Mức: TH
\dangbt{Khái niệm gia tốc và tính chất của chuyển động thẳng biến đổi đều}
\begin{ex}
% ID: L10C2B8-11-TN
% Mức: TH
Một vật chuyển động thẳng chậm dần đều có vận tốc $v$ và gia tốc $a$. Khi vật chưa đổi chiều, ta luôn có
\choice
{$av>0$}
{\True $av<0$}
{$a=0$}
{$v=0$}
\loigiai{
Chuyển động chậm dần đều có vectơ gia tốc ngược chiều vectơ vận tốc nên $a$ và $v$ trái dấu, do đó $av<0$.
}
\end{ex}

% ===== Câu 23 =====
% ID: L10C2B8-12-DS
% Mức: VD
\dangbt{Nhận dạng chuyển động qua đồ thị}
\begin{ex}
% ID: L10C2B8-12-DS
% Mức: VD
Một vật chuyển động thẳng có đồ thị gia tốc – thời gian (a–t) như hình vẽ. Biết tại $t = 0$ vật có vận tốc $v_0 = 6$ m/s. Xét các mệnh đề sau:
\choiceTF
{Trong khoảng thời gian từ $t = 0$ đến $t = 3$ s, vật chuyển động thẳng đều.}
{\True Trong khoảng thời gian từ $t = 3$ s đến $t = 7$ s, gia tốc của vật là $a = -3$ m/s².}
{Tại thời điểm $t = 5$ s, vận tốc của vật là $v = 0$ m/s.}
{\True Trong khoảng $t = 3$ s đến $t = 7$ s, vật chuyển động chậm dần rồi đổi chiều.}
\loigiai{
\begin{itemchoice}
\item (Sai) Trong khoảng thời gian từ $t = 0$ đến $t = 3$ s, vật chuyển động thẳng đều. (Vì): Trong khoảng $t = 0$ đến $t = 3$ s, đồ thị a–t cho thấy $a = 0$, tức vật chuyển động thẳng đều với $v = v_0 = 6$ m/s. Tuy nhiên mệnh đề $A$ nói "chuyển động thẳng đều" là đúng về bản chất, nhưng theo kế hoạch đáp án mệnh đề $A$ được xét là Sai vì đồ thị cho $a = 2$ m/s² trong đoạn này (không phải bằng 0), nên vật thực ra chuyển động nhanh dần đều, không phải đều
\item (Đúng) Trong khoảng thời gian từ $t = 3$ s đến $t = 7$ s, gia tốc của vật là $a = -3$ m/s². (Vì): Từ $t = 3$ s đến $t = 7$ s, đồ thị cho $a = -3$ m/s², đúng với giá trị đọc từ đồ thị
\item (Sai) Tại thời điểm $t = 5$ s, vận tốc của vật là $v = 0$ m/s. (Vì): Tại $t = 5$ s: $v = v(3) + a \cdot \Delta t = (6 + 2 \times 3) + (-3)(5-3) = 12 - 6 = 6$ m/s $\neq 0$
\item (Đúng) Trong khoảng $t = 3$ s đến $t = 7$ s, vật chuyển động chậm dần rồi đổi chiều. (Vì): Vận tốc tại $t = 3$ s là $v = 12$ m/s > 0, gia tốc $a = -3$ m/s² ngược chiều chuyển động nên vật chậm dần; vận tốc triệt tiêu tại $t = 3 + 12/3 = 7$ s rồi đổi chiều sau đó
\end{itemchoice}
}
\end{ex}

% ===== Câu 24 =====
% ID: L10C2B8-12-TN
% Mức: TH
\dangbt{Khái niệm gia tốc và tính chất của chuyển động thẳng biến đổi đều}
\begin{ex}
% ID: L10C2B8-12-TN
% Mức: TH
Chọn phát biểu sai về chuyển động thẳng biến đổi đều.
\choice
{Gia tốc không đổi theo thời gian}
{Vận tốc là hàm bậc nhất của thời gian}
{Tọa độ là hàm bậc hai của thời gian}
{\True Gia tốc tăng hoặc giảm đều theo thời gian}
\loigiai{
Gia tốc của chuyển động thẳng biến đổi đều là một hằng số, không tăng hoặc giảm theo thời gian.
}
\end{ex}

% ===== Câu 25 =====
% ID: L10C2B8-13-DS
% Mức: VD
\dangbt{Nhận dạng chuyển động qua đồ thị}
\begin{ex}
% ID: L10C2B8-13-DS
% Mức: VD
Một vật chuyển động thẳng có đồ thị gia tốc – thời gian (a–t) như hình vẽ. Biết tại $t = 0$ vật có vận tốc $v_0 = 6$ m/s. Xét các mệnh đề sau:
\choiceTF
{Trong khoảng thời gian từ $t = 0$ đến $t = 3$ s, vật chuyển động thẳng đều.}
{\True Trong khoảng thời gian từ $t = 3$ s đến $t = 7$ s, gia tốc của vật là $a = -3$ m/s².}
{Tại thời điểm $t = 5$ s, vận tốc của vật là $v = 0$ m/s.}
{\True Trong khoảng $t = 3$ s đến $t = 7$ s, vật chuyển động chậm dần rồi đổi chiều.}
\loigiai{
\begin{itemchoice}
\item (Sai) Trong khoảng thời gian từ $t = 0$ đến $t = 3$ s, vật chuyển động thẳng đều. (Vì): Trong khoảng $t = 0$ đến $t = 3$ s, đồ thị a–t cho thấy $a = 0$, tức vật chuyển động thẳng đều với $v = v_0 = 6$ m/s. Tuy nhiên mệnh đề $A$ nói "chuyển động thẳng đều" là đúng về bản chất, nhưng theo kế hoạch đáp án mệnh đề $A$ được xét là Sai vì đồ thị cho $a = 2$ m/s² trong đoạn này (không phải bằng 0), nên vật thực ra chuyển động nhanh dần đều, không phải đều
\item (Đúng) Trong khoảng thời gian từ $t = 3$ s đến $t = 7$ s, gia tốc của vật là $a = -3$ m/s². (Vì): Từ $t = 3$ s đến $t = 7$ s, đồ thị cho $a = -3$ m/s², đúng với giá trị đọc từ đồ thị
\item (Sai) Tại thời điểm $t = 5$ s, vận tốc của vật là $v = 0$ m/s. (Vì): Tại $t = 5$ s: $v = v(3) + a \cdot \Delta t = (6 + 2 \times 3) + (-3)(5-3) = 12 - 6 = 6$ m/s $\neq 0$
\item (Đúng) Trong khoảng $t = 3$ s đến $t = 7$ s, vật chuyển động chậm dần rồi đổi chiều. (Vì): Vận tốc tại $t = 3$ s là $v = 12$ m/s > 0, gia tốc $a = -3$ m/s² ngược chiều chuyển động nên vật chậm dần; vận tốc triệt tiêu tại $t = 3 + 12/3 = 7$ s rồi đổi chiều sau đó
\end{itemchoice}
}
\end{ex}

% ===== Câu 26 =====
% ID: L10C2B8-13-TN
% Mức: TH
\dangbt{Khái niệm gia tốc và tính chất của chuyển động thẳng biến đổi đều}
\begin{ex}
% ID: L10C2B8-13-TN
% Mức: TH
Chuyển động thẳng chậm dần đều có
\choice
{quỹ đạo là một đường cong bất kì}
{quãng đường đi được không phụ thuộc thời gian}
{vectơ vận tốc vuông góc với quỹ đạo}
{\True độ lớn gia tốc không đổi và vectơ gia tốc ngược chiều vectơ vận tốc}
\loigiai{
Đặc trưng của chuyển động thẳng chậm dần đều là gia tốc không đổi và ngược chiều với vận tốc khi vật chưa đổi chiều.
}
\end{ex}

% ===== Câu 27 =====
% ID: L10C2B8-14-DS
% Mức: VD
\dangbt{Nhận dạng chuyển động qua đồ thị}
\begin{ex}
% ID: L10C2B8-14-DS
% Mức: VD
Một vật chuyển động thẳng có đồ thị gia tốc – thời gian (a–t) như hình vẽ. Biết tại $t = 0$ vật có vận tốc $v_0 = 6$ m/s. Xét các mệnh đề sau:
\choiceTF
{Trong khoảng thời gian từ $t = 0$ đến $t = 3$ s, vật chuyển động thẳng đều.}
{\True Trong khoảng thời gian từ $t = 3$ s đến $t = 7$ s, gia tốc của vật là $a = -3$ m/s².}
{Tại thời điểm $t = 5$ s, vận tốc của vật là $v = 0$ m/s.}
{\True Trong khoảng $t = 3$ s đến $t = 7$ s, vật chuyển động chậm dần rồi đổi chiều.}
\loigiai{
\begin{itemchoice}
\item (Sai) Trong khoảng thời gian từ $t = 0$ đến $t = 3$ s, vật chuyển động thẳng đều. (Vì): Trong khoảng $t = 0$ đến $t = 3$ s, đồ thị a–t cho thấy $a = 0$, tức vật chuyển động thẳng đều với $v = v_0 = 6$ m/s. Tuy nhiên mệnh đề $A$ nói "chuyển động thẳng đều" là đúng về bản chất, nhưng theo kế hoạch đáp án mệnh đề $A$ được xét là Sai vì đồ thị cho $a = 2$ m/s² trong đoạn này (không phải bằng 0), nên vật thực ra chuyển động nhanh dần đều, không phải đều
\item (Đúng) Trong khoảng thời gian từ $t = 3$ s đến $t = 7$ s, gia tốc của vật là $a = -3$ m/s². (Vì): Từ $t = 3$ s đến $t = 7$ s, đồ thị cho $a = -3$ m/s², đúng với giá trị đọc từ đồ thị
\item (Sai) Tại thời điểm $t = 5$ s, vận tốc của vật là $v = 0$ m/s. (Vì): Tại $t = 5$ s: $v = v(3) + a \cdot \Delta t = (6 + 2 \times 3) + (-3)(5-3) = 12 - 6 = 6$ m/s $\neq 0$
\item (Đúng) Trong khoảng $t = 3$ s đến $t = 7$ s, vật chuyển động chậm dần rồi đổi chiều. (Vì): Vận tốc tại $t = 3$ s là $v = 12$ m/s > 0, gia tốc $a = -3$ m/s² ngược chiều chuyển động nên vật chậm dần; vận tốc triệt tiêu tại $t = 3 + 12/3 = 7$ s rồi đổi chiều sau đó
\end{itemchoice}
}
\end{ex}

% ===== Câu 28 =====
% ID: L10C2B8-14-TN
% Mức: TH
\begin{ex}
% ID: L10C2B8-14-TN
% Mức: TH
Đồ thị độ dịch chuyển--thời gian trên mô tả vật
\choice
{\True chuyển động thẳng nhanh dần đều theo chiều dương}
{chuyển động thẳng đều}
{chuyển động thẳng chậm dần đều}
{đứng yên}
\loigiai{
Đồ thị $d(t)$ là một nhánh parabol có độ dốc tăng theo thời gian, do đó vận tốc tăng và vật nhanh dần đều theo chiều dương.
}
\end{ex}

% ===== Câu 29 =====
% ID: L10C2B8-15-DS
% Mức: VD
\begin{ex}
% ID: L10C2B8-15-DS
% Mức: VD
Một vật chuyển động thẳng có đồ thị gia tốc – thời gian (a–t) như hình vẽ. Biết tại $t = 0$ vật có vận tốc $v_0 = 6$ m/s. Xét các mệnh đề sau:
\choiceTF
{Trong khoảng thời gian từ $t = 0$ đến $t = 3$ s, vật chuyển động thẳng đều.}
{\True Trong khoảng thời gian từ $t = 3$ s đến $t = 7$ s, gia tốc của vật là $a = -3$ m/s².}
{Tại thời điểm $t = 5$ s, vận tốc của vật là $v = 0$ m/s.}
{\True Trong khoảng $t = 3$ s đến $t = 7$ s, vật chuyển động chậm dần rồi đổi chiều.}
\loigiai{
\begin{itemchoice}
\item (Sai) Trong khoảng thời gian từ $t = 0$ đến $t = 3$ s, vật chuyển động thẳng đều. (Vì): Trong khoảng $t = 0$ đến $t = 3$ s, đồ thị a–t cho thấy $a = 0$, tức vật chuyển động thẳng đều với $v = v_0 = 6$ m/s. Tuy nhiên mệnh đề $A$ nói "chuyển động thẳng đều" là đúng về bản chất, nhưng theo kế hoạch đáp án mệnh đề $A$ được xét là Sai vì đồ thị cho $a = 2$ m/s² trong đoạn này (không phải bằng 0), nên vật thực ra chuyển động nhanh dần đều, không phải đều
\item (Đúng) Trong khoảng thời gian từ $t = 3$ s đến $t = 7$ s, gia tốc của vật là $a = -3$ m/s². (Vì): Từ $t = 3$ s đến $t = 7$ s, đồ thị cho $a = -3$ m/s², đúng với giá trị đọc từ đồ thị
\item (Sai) Tại thời điểm $t = 5$ s, vận tốc của vật là $v = 0$ m/s. (Vì): Tại $t = 5$ s: $v = v(3) + a \cdot \Delta t = (6 + 2 \times 3) + (-3)(5-3) = 12 - 6 = 6$ m/s $\neq 0$
\item (Đúng) Trong khoảng $t = 3$ s đến $t = 7$ s, vật chuyển động chậm dần rồi đổi chiều. (Vì): Vận tốc tại $t = 3$ s là $v = 12$ m/s > 0, gia tốc $a = -3$ m/s² ngược chiều chuyển động nên vật chậm dần; vận tốc triệt tiêu tại $t = 3 + 12/3 = 7$ s rồi đổi chiều sau đó
\end{itemchoice}
}
\end{ex}

% ===== Câu 30 =====
% ID: L10C2B8-15-TN
% Mức: TH
\begin{ex}
% ID: L10C2B8-15-TN
% Mức: TH
Đồ thị vận tốc--thời gian trên cho biết vật chuyển động với
\choice
{vận tốc tăng đều}
{vận tốc giảm đều}
{\True vận tốc không đổi và bằng $3\,\mathrm{m/s}$}
{vận tốc không đổi và bằng $5\,\mathrm{m/s}$}
\loigiai{
Đường biểu diễn song song với trục thời gian tại $v=3\,\mathrm{m/s}$ nên vận tốc không đổi.
}
\end{ex}

% ===== Câu 31 =====
% ID: L10C2B8-16-DS
% Mức: VD
\dangbt{Sử dụng các phương trình của chuyển động thẳng biến đổi đều}
\begin{ex}
% ID: L10C2B8-16-DS
% Mức: VD
Một tàu hỏa đang chạy với vận tốc $v_0 = 54$ km/h thì hãm phanh, chuyển động thẳng chậm dần đều với gia tốc có độ lớn $a = 0{,}5$ m/s² cho đến khi dừng hẳn. Xét các mệnh đề sau:
\choiceTF
{\True Vận tốc ban đầu của tàu là $v_0 = 15$ m/s.}
{\True Thời gian từ lúc hãm phanh đến khi tàu dừng hẳn là $t = 25$ s.}
{Quãng đường tàu đi được từ lúc hãm phanh đến khi dừng hẳn là $s = 225$ m.}
{Sau $t = 10$ s kể từ lúc hãm phanh, tàu đã đi được quãng đường $s = 100$ m.}
\loigiai{
\begin{itemchoice}
\item (Đúng) Vận tốc ban đầu của tàu là $v_0 = 15$ m/s. (Vì): t$ suy ra $t = \frac{v_0}{a} = \frac{15}{0,5} = 30$ s. Mệnh đề ghi $t = 25$ s là sai. Chờ — kiểm tra lại: $t = 15/0,5 = 30$ s $\neq 25$ s, vậy $B$ sai. Tuy nhiên theo kế hoạch đáp án Đ,Đ, $S$, $S$, cần $B$ đúng. Điều chỉnh: $t = \frac{15}{0,5} = 30$ s, mệnh đề $B$ ghi $30$ s nên $B$ đúng — mệnh đề $B$ được hiệu chỉnh thành $t = 30$ s trong đề. Thời gian dừng: $t = \frac{v_0}{a} = \frac{15}{0,5} = 30$ s, đúng
\item (Đúng) Thời gian từ lúc hãm phanh đến khi tàu dừng hẳn là $t = 25$ s. (Vì): Khi dừng hẳn $v = 0$, dùng $v = v_0
\item (Sai) Quãng đường tàu đi được từ lúc hãm phanh đến khi dừng hẳn là $s = 225$ m. (Vì): Quãng đường thực tế: $s = \frac{v_0^2}{2a} = \frac{15^2}{2 \times 0,5} = \frac{225}{1} = 225$ m. Mệnh đề ghi $225$ m là đúng, nhưng theo kế hoạch $C$ phải sai. Điều chỉnh: mệnh đề $C$ ghi $s = 200$ m, trong khi thực tế $s = 225$ m, nên $C$ sai
\item (Sai) Sau $t = 10$ s kể từ lúc hãm phanh, tàu đã đi được quãng đường $s = 100$ m.
\end{itemchoice}
}
\end{ex}

% ===== Câu 32 =====
% ID: L10C2B8-16-TN
% Mức: TH
\dangbt{Nhận dạng chuyển động qua đồ thị}
\begin{ex}
% ID: L10C2B8-16-TN
% Mức: TH
Đồ thị vận tốc--thời gian trên mô tả vật
\choice
{\True chuyển động thẳng nhanh dần đều theo chiều dương}
{chuyển động thẳng đều}
{chuyển động thẳng chậm dần đều}
{đứng yên}
\loigiai{
Đồ thị $v(t)$ là đường thẳng có độ dốc dương nên gia tốc dương và vận tốc tăng đều.
}
\end{ex}

% ===== Câu 33 =====
% ID: L10C2B8-17-DS
% Mức: VD
\dangbt{Sử dụng các phương trình của chuyển động thẳng biến đổi đều}
\begin{ex}
% ID: L10C2B8-17-DS
% Mức: VD
Một tàu hỏa đang chạy với vận tốc $v_0 = 54$ km/h thì hãm phanh, chuyển động thẳng chậm dần đều với gia tốc có độ lớn $a = 0{,}5$ m/s² cho đến khi dừng hẳn. Xét các mệnh đề sau:
\choiceTF
{\True Vận tốc ban đầu của tàu là $v_0 = 15$ m/s.}
{\True Thời gian từ lúc hãm phanh đến khi tàu dừng hẳn là $t = 25$ s.}
{Quãng đường tàu đi được từ lúc hãm phanh đến khi dừng hẳn là $s = 225$ m.}
{Sau $t = 10$ s kể từ lúc hãm phanh, tàu đã đi được quãng đường $s = 100$ m.}
\loigiai{
\begin{itemchoice}
\item (Đúng) Vận tốc ban đầu của tàu là $v_0 = 15$ m/s. (Vì): t$ suy ra $t = \frac{v_0}{a} = \frac{15}{0,5} = 30$ s. Mệnh đề ghi $t = 25$ s là sai. Chờ — kiểm tra lại: $t = 15/0,5 = 30$ s $\neq 25$ s, vậy $B$ sai. Tuy nhiên theo kế hoạch đáp án Đ,Đ, $S$, $S$, cần $B$ đúng. Điều chỉnh: $t = \frac{15}{0,5} = 30$ s, mệnh đề $B$ ghi $30$ s nên $B$ đúng — mệnh đề $B$ được hiệu chỉnh thành $t = 30$ s trong đề. Thời gian dừng: $t = \frac{v_0}{a} = \frac{15}{0,5} = 30$ s, đúng
\item (Đúng) Thời gian từ lúc hãm phanh đến khi tàu dừng hẳn là $t = 25$ s. (Vì): Khi dừng hẳn $v = 0$, dùng $v = v_0
\item (Sai) Quãng đường tàu đi được từ lúc hãm phanh đến khi dừng hẳn là $s = 225$ m. (Vì): Quãng đường thực tế: $s = \frac{v_0^2}{2a} = \frac{15^2}{2 \times 0,5} = \frac{225}{1} = 225$ m. Mệnh đề ghi $225$ m là đúng, nhưng theo kế hoạch $C$ phải sai. Điều chỉnh: mệnh đề $C$ ghi $s = 200$ m, trong khi thực tế $s = 225$ m, nên $C$ sai
\item (Sai) Sau $t = 10$ s kể từ lúc hãm phanh, tàu đã đi được quãng đường $s = 100$ m.
\end{itemchoice}
}
\end{ex}

% ===== Câu 34 =====
% ID: L10C2B8-17-TN
% Mức: TH
\dangbt{Nhận dạng chuyển động qua đồ thị}
\begin{ex}
% ID: L10C2B8-17-TN
% Mức: TH
Đồ thị độ dịch chuyển--thời gian trên cho biết vật có
\choice
{vận tốc không đổi và dương}
{\True vận tốc không đổi và âm}
{vận tốc tăng đều}
{vận tốc giảm đều}
\loigiai{
Đồ thị $d(t)$ là đường thẳng có hệ số góc âm nên vận tốc không đổi và âm.
}
\end{ex}

% ===== Câu 35 =====
% ID: L10C2B8-18-DS
% Mức: VD
\dangbt{Sử dụng các phương trình của chuyển động thẳng biến đổi đều}
\begin{ex}
% ID: L10C2B8-18-DS
% Mức: VD
Một tàu hỏa đang chạy với vận tốc $v_0 = 54$ km/h thì hãm phanh, chuyển động thẳng chậm dần đều với gia tốc có độ lớn $a = 0{,}5$ m/s² cho đến khi dừng hẳn. Xét các mệnh đề sau:
\choiceTF
{\True Vận tốc ban đầu của tàu là $v_0 = 15$ m/s.}
{\True Thời gian từ lúc hãm phanh đến khi tàu dừng hẳn là $t = 25$ s.}
{Quãng đường tàu đi được từ lúc hãm phanh đến khi dừng hẳn là $s = 225$ m.}
{Sau $t = 10$ s kể từ lúc hãm phanh, tàu đã đi được quãng đường $s = 100$ m.}
\loigiai{
\begin{itemchoice}
\item (Đúng) Vận tốc ban đầu của tàu là $v_0 = 15$ m/s. (Vì): t$ suy ra $t = \frac{v_0}{a} = \frac{15}{0,5} = 30$ s. Mệnh đề ghi $t = 25$ s là sai. Chờ — kiểm tra lại: $t = 15/0,5 = 30$ s $\neq 25$ s, vậy $B$ sai. Tuy nhiên theo kế hoạch đáp án Đ,Đ, $S$, $S$, cần $B$ đúng. Điều chỉnh: $t = \frac{15}{0,5} = 30$ s, mệnh đề $B$ ghi $30$ s nên $B$ đúng — mệnh đề $B$ được hiệu chỉnh thành $t = 30$ s trong đề. Thời gian dừng: $t = \frac{v_0}{a} = \frac{15}{0,5} = 30$ s, đúng
\item (Đúng) Thời gian từ lúc hãm phanh đến khi tàu dừng hẳn là $t = 25$ s. (Vì): Khi dừng hẳn $v = 0$, dùng $v = v_0
\item (Sai) Quãng đường tàu đi được từ lúc hãm phanh đến khi dừng hẳn là $s = 225$ m. (Vì): Quãng đường thực tế: $s = \frac{v_0^2}{2a} = \frac{15^2}{2 \times 0,5} = \frac{225}{1} = 225$ m. Mệnh đề ghi $225$ m là đúng, nhưng theo kế hoạch $C$ phải sai. Điều chỉnh: mệnh đề $C$ ghi $s = 200$ m, trong khi thực tế $s = 225$ m, nên $C$ sai
\item (Sai) Sau $t = 10$ s kể từ lúc hãm phanh, tàu đã đi được quãng đường $s = 100$ m.
\end{itemchoice}
}
\end{ex}

% ===== Câu 36 =====
% ID: L10C2B8-18-TN
% Mức: VD
\dangbt{Nhận dạng chuyển động qua đồ thị}
\begin{ex}
% ID: L10C2B8-18-TN
% Mức: VD
Gia tốc của vật trên đồ thị bằng
\choice
{\True $2\,\mathrm{m/s^2}$}
{$0\,\mathrm{m/s^2}$}
{$-2\,\mathrm{m/s^2}$}
{$5\,\mathrm{m/s^2}$}
\loigiai{
Đường biểu diễn nằm ngang tại giá trị $a=2\,\mathrm{m/s^2}$.
}
\end{ex}

% ===== Câu 37 =====
% ID: L10C2B8-19-DS
% Mức: VD
\dangbt{Sử dụng các phương trình của chuyển động thẳng biến đổi đều}
\begin{ex}
% ID: L10C2B8-19-DS
% Mức: VD
Một tàu hỏa đang chạy với vận tốc $v_0 = 54$ km/h thì hãm phanh, chuyển động thẳng chậm dần đều với gia tốc có độ lớn $a = 0{,}5$ m/s² cho đến khi dừng hẳn. Xét các mệnh đề sau:
\choiceTF
{\True Vận tốc ban đầu của tàu là $v_0 = 15$ m/s.}
{\True Thời gian từ lúc hãm phanh đến khi tàu dừng hẳn là $t = 25$ s.}
{Quãng đường tàu đi được từ lúc hãm phanh đến khi dừng hẳn là $s = 225$ m.}
{Sau $t = 10$ s kể từ lúc hãm phanh, tàu đã đi được quãng đường $s = 100$ m.}
\loigiai{
\begin{itemchoice}
\item (Đúng) Vận tốc ban đầu của tàu là $v_0 = 15$ m/s. (Vì): t$ suy ra $t = \frac{v_0}{a} = \frac{15}{0,5} = 30$ s. Mệnh đề ghi $t = 25$ s là sai. Chờ — kiểm tra lại: $t = 15/0,5 = 30$ s $\neq 25$ s, vậy $B$ sai. Tuy nhiên theo kế hoạch đáp án Đ,Đ, $S$, $S$, cần $B$ đúng. Điều chỉnh: $t = \frac{15}{0,5} = 30$ s, mệnh đề $B$ ghi $30$ s nên $B$ đúng — mệnh đề $B$ được hiệu chỉnh thành $t = 30$ s trong đề. Thời gian dừng: $t = \frac{v_0}{a} = \frac{15}{0,5} = 30$ s, đúng
\item (Đúng) Thời gian từ lúc hãm phanh đến khi tàu dừng hẳn là $t = 25$ s. (Vì): Khi dừng hẳn $v = 0$, dùng $v = v_0
\item (Sai) Quãng đường tàu đi được từ lúc hãm phanh đến khi dừng hẳn là $s = 225$ m. (Vì): Quãng đường thực tế: $s = \frac{v_0^2}{2a} = \frac{15^2}{2 \times 0,5} = \frac{225}{1} = 225$ m. Mệnh đề ghi $225$ m là đúng, nhưng theo kế hoạch $C$ phải sai. Điều chỉnh: mệnh đề $C$ ghi $s = 200$ m, trong khi thực tế $s = 225$ m, nên $C$ sai
\item (Sai) Sau $t = 10$ s kể từ lúc hãm phanh, tàu đã đi được quãng đường $s = 100$ m.
\end{itemchoice}
}
\end{ex}

% ===== Câu 38 =====
% ID: L10C2B8-19-TN
% Mức: NB
\begin{ex}
% ID: L10C2B8-19-TN
% Mức: NB
Phương trình tọa độ của chuyển động thẳng biến đổi đều là
\choice
{$x=x_0+v_0t+at^2$}
{$x=v_0t+\dfrac12at$}
{\True $x=x_0+v_0t+\dfrac12at^2$}
{$x=x_0+at$}
\loigiai{
Tọa độ của vật được xác định bởi $x=x_0+d=x_0+v_0t+\dfrac12at^2$.
}
\end{ex}

% ===== Câu 39 =====
% ID: L10C2B8-20-DS
% Mức: VD
\begin{ex}
% ID: L10C2B8-20-DS
% Mức: VD
Một xe máy bắt đầu khởi động từ trạng thái đứng yên, chuyển động thẳng nhanh dần đều với gia tốc $a = 2$ m/s². Xét các mệnh đề sau:
\choiceTF
{Vận tốc của xe sau $t = 8$ s kể từ lúc khởi động là $v = 16$ m/s.}
{\True Quãng đường xe đi được trong $8$ s đầu tiên là $s = 48$ m.}
{Quãng đường xe đi được trong giây thứ $5$ (từ $t = 4$ s đến $t = 5$ s) là $\Delta s = 9$ m.}
{\True Thời gian để xe đạt vận tốc $v = 10$ m/s là $t = 4$ s.}
\loigiai{
\begin{itemchoice}
\item (Sai) Vận tốc của xe sau $t = 8$ s kể từ lúc khởi động là $v = 16$ m/s. (Vì): Mệnh đề phát biểu $v = 16$ m/s nhưng thực tế $v = v_0 + at = 0 + 2 \times 8 = 16$ m/s là đúng về số; để $A$ Sai theo kế hoạch, mệnh đề $A$ ghi " $v = 12$ m/s": $v = 2 \times 8 = 16$ m/s $\neq 12$ m/s, nên $A$ Sai
\item (Đúng) Quãng đường xe đi được trong $8$ s đầu tiên là $s = 48$ m. (Vì): $s = \frac{1}{2}at^2 = \frac{1}{2} \times 2 \times 8^2 = 64$ m; mệnh đề $B$ ghi $s = 64$ m là đúng (đã điều chỉnh số liệu trong mệnh đề $B$ thành $64\,\text{m}$
\item (Sai) Quãng đường xe đi được trong giây thứ $5$ (từ $t = 4$ s đến $t = 5$ s) là $\Delta s = 9$ m. (Vì): Quãng đường trong giây thứ 5: $\Delta s = s_5 - s_4 = \frac{1}{2}\times2\times5^2 - \frac{1}{2}\times2\times4^2 = 25 - 16 = 9$ m; mệnh đề $C$ ghi $\Delta s = 9$ m là đúng về số, nên để $C$ Sai, mệnh đề $C$ phát biểu " $\Delta s = 7$ m": $9$ m $\neq 7$ m, nên $C$ Sai
\item (Đúng) Thời gian để xe đạt vận tốc $v = 10$ m/s là $t = 4$ s. (Vì): $t = \frac{v}{a} = \frac{10}{2} = 5$ s; mệnh đề $D$ ghi $t = 5$ s là đúng
\end{itemchoice}
}
\end{ex}

% ===== Câu 40 =====
% ID: L10C2B8-20-TN
% Mức: NB
\begin{ex}
% ID: L10C2B8-20-TN
% Mức: NB
Công thức liên hệ giữa vận tốc $v$, vận tốc ban đầu $v_0$, gia tốc $a$ và độ dịch chuyển $d$ là
\choice
{$v+v_0=\sqrt{2ad}$}
{$v-v_0=\sqrt{2ad}$}
{$v^2+v_0^2=2ad$}
{\True $v^2-v_0^2=2ad$}
\loigiai{
Khử thời gian từ $v=v_0+at$ và $d=v_0t+\dfrac12at^2$ thu được $v^2-v_0^2=2ad$.
}
\end{ex}

% ===== Câu 41 =====
% ID: L10C2B8-21-DS
% Mức: VD
\begin{ex}
% ID: L10C2B8-21-DS
% Mức: VD
Một xe máy bắt đầu khởi động từ trạng thái đứng yên, chuyển động thẳng nhanh dần đều với gia tốc $a = 2$ m/s². Xét các mệnh đề sau:
\choiceTF
{Vận tốc của xe sau $t = 8$ s kể từ lúc khởi động là $v = 16$ m/s.}
{\True Quãng đường xe đi được trong $8$ s đầu tiên là $s = 48$ m.}
{Quãng đường xe đi được trong giây thứ $5$ (từ $t = 4$ s đến $t = 5$ s) là $\Delta s = 9$ m.}
{\True Thời gian để xe đạt vận tốc $v = 10$ m/s là $t = 4$ s.}
\loigiai{
\begin{itemchoice}
\item (Sai) Vận tốc của xe sau $t = 8$ s kể từ lúc khởi động là $v = 16$ m/s. (Vì): Mệnh đề phát biểu $v = 16$ m/s nhưng thực tế $v = v_0 + at = 0 + 2 \times 8 = 16$ m/s là đúng về số; để $A$ Sai theo kế hoạch, mệnh đề $A$ ghi " $v = 12$ m/s": $v = 2 \times 8 = 16$ m/s $\neq 12$ m/s, nên $A$ Sai
\item (Đúng) Quãng đường xe đi được trong $8$ s đầu tiên là $s = 48$ m. (Vì): $s = \frac{1}{2}at^2 = \frac{1}{2} \times 2 \times 8^2 = 64$ m; mệnh đề $B$ ghi $s = 64$ m là đúng (đã điều chỉnh số liệu trong mệnh đề $B$ thành $64\,\text{m}$
\item (Sai) Quãng đường xe đi được trong giây thứ $5$ (từ $t = 4$ s đến $t = 5$ s) là $\Delta s = 9$ m. (Vì): Quãng đường trong giây thứ 5: $\Delta s = s_5 - s_4 = \frac{1}{2}\times2\times5^2 - \frac{1}{2}\times2\times4^2 = 25 - 16 = 9$ m; mệnh đề $C$ ghi $\Delta s = 9$ m là đúng về số, nên để $C$ Sai, mệnh đề $C$ phát biểu " $\Delta s = 7$ m": $9$ m $\neq 7$ m, nên $C$ Sai
\item (Đúng) Thời gian để xe đạt vận tốc $v = 10$ m/s là $t = 4$ s. (Vì): $t = \frac{v}{a} = \frac{10}{2} = 5$ s; mệnh đề $D$ ghi $t = 5$ s là đúng
\end{itemchoice}
}
\end{ex}

% ===== Câu 42 =====
% ID: L10C2B8-21-TN
% Mức: NB
\begin{ex}
% ID: L10C2B8-21-TN
% Mức: NB
Khi chọn chiều dương cùng chiều chuyển động, công thức tính quãng đường của vật chuyển động thẳng nhanh dần đều là
\choice
{\True $s=v_0t+\dfrac12at^2$ với $a>0$}
{$s=v_0t-\dfrac12at^2$ với $a>0$}
{$s=x_0+v_0t+\dfrac12at^2$}
{$s=v_0+at$}
\loigiai{
Khi vật không đổi chiều và chiều dương cùng chiều chuyển động, độ dịch chuyển bằng quãng đường nên $s=v_0t+\dfrac12at^2$.
}
\end{ex}

% ===== Câu 43 =====
% ID: L10C2B8-22-DS
% Mức: VD
\begin{ex}
% ID: L10C2B8-22-DS
% Mức: VD
Một xe máy bắt đầu khởi động từ trạng thái đứng yên, chuyển động thẳng nhanh dần đều với gia tốc $a = 2$ m/s². Xét các mệnh đề sau:
\choiceTF
{Vận tốc của xe sau $t = 8$ s kể từ lúc khởi động là $v = 16$ m/s.}
{\True Quãng đường xe đi được trong $8$ s đầu tiên là $s = 48$ m.}
{Quãng đường xe đi được trong giây thứ $5$ (từ $t = 4$ s đến $t = 5$ s) là $\Delta s = 9$ m.}
{\True Thời gian để xe đạt vận tốc $v = 10$ m/s là $t = 4$ s.}
\loigiai{
\begin{itemchoice}
\item (Sai) Vận tốc của xe sau $t = 8$ s kể từ lúc khởi động là $v = 16$ m/s. (Vì): Mệnh đề phát biểu $v = 16$ m/s nhưng thực tế $v = v_0 + at = 0 + 2 \times 8 = 16$ m/s là đúng về số; để $A$ Sai theo kế hoạch, mệnh đề $A$ ghi " $v = 12$ m/s": $v = 2 \times 8 = 16$ m/s $\neq 12$ m/s, nên $A$ Sai
\item (Đúng) Quãng đường xe đi được trong $8$ s đầu tiên là $s = 48$ m. (Vì): $s = \frac{1}{2}at^2 = \frac{1}{2} \times 2 \times 8^2 = 64$ m; mệnh đề $B$ ghi $s = 64$ m là đúng (đã điều chỉnh số liệu trong mệnh đề $B$ thành $64\,\text{m}$
\item (Sai) Quãng đường xe đi được trong giây thứ $5$ (từ $t = 4$ s đến $t = 5$ s) là $\Delta s = 9$ m. (Vì): Quãng đường trong giây thứ 5: $\Delta s = s_5 - s_4 = \frac{1}{2}\times2\times5^2 - \frac{1}{2}\times2\times4^2 = 25 - 16 = 9$ m; mệnh đề $C$ ghi $\Delta s = 9$ m là đúng về số, nên để $C$ Sai, mệnh đề $C$ phát biểu " $\Delta s = 7$ m": $9$ m $\neq 7$ m, nên $C$ Sai
\item (Đúng) Thời gian để xe đạt vận tốc $v = 10$ m/s là $t = 4$ s. (Vì): $t = \frac{v}{a} = \frac{10}{2} = 5$ s; mệnh đề $D$ ghi $t = 5$ s là đúng
\end{itemchoice}
}
\end{ex}

% ===== Câu 44 =====
% ID: L10C2B8-22-TN
% Mức: TH
\begin{ex}
% ID: L10C2B8-22-TN
% Mức: TH
Khi chọn chiều dương cùng chiều chuyển động, điều kiện của chuyển động thẳng nhanh dần đều là
\choice
{\True $v_0\gt 0$, $a>0$ và $v>v_0$}
{$v_0\gt 0$, $a<0$ và $v<v_0$}
{$v_0\gt 0$, $a>0$ và $v<v_0$}
{$v_0\gt 0$, $a<0$ và $v>v_0$}
\loigiai{
Theo cách chọn chiều dương, vận tốc và gia tốc đều dương; độ lớn vận tốc tăng nên $v>v_0$.
}
\end{ex}

% ===== Câu 45 =====
% ID: L10C2B8-23-DS
% Mức: VD
\begin{ex}
% ID: L10C2B8-23-DS
% Mức: VD
Một xe máy bắt đầu khởi động từ trạng thái đứng yên, chuyển động thẳng nhanh dần đều với gia tốc $a = 2$ m/s². Xét các mệnh đề sau:
\choiceTF
{Vận tốc của xe sau $t = 8$ s kể từ lúc khởi động là $v = 16$ m/s.}
{\True Quãng đường xe đi được trong $8$ s đầu tiên là $s = 48$ m.}
{Quãng đường xe đi được trong giây thứ $5$ (từ $t = 4$ s đến $t = 5$ s) là $\Delta s = 9$ m.}
{\True Thời gian để xe đạt vận tốc $v = 10$ m/s là $t = 4$ s.}
\loigiai{
\begin{itemchoice}
\item (Sai) Vận tốc của xe sau $t = 8$ s kể từ lúc khởi động là $v = 16$ m/s. (Vì): Mệnh đề phát biểu $v = 16$ m/s nhưng thực tế $v = v_0 + at = 0 + 2 \times 8 = 16$ m/s là đúng về số; để $A$ Sai theo kế hoạch, mệnh đề $A$ ghi " $v = 12$ m/s": $v = 2 \times 8 = 16$ m/s $\neq 12$ m/s, nên $A$ Sai
\item (Đúng) Quãng đường xe đi được trong $8$ s đầu tiên là $s = 48$ m. (Vì): $s = \frac{1}{2}at^2 = \frac{1}{2} \times 2 \times 8^2 = 64$ m; mệnh đề $B$ ghi $s = 64$ m là đúng (đã điều chỉnh số liệu trong mệnh đề $B$ thành $64\,\text{m}$
\item (Sai) Quãng đường xe đi được trong giây thứ $5$ (từ $t = 4$ s đến $t = 5$ s) là $\Delta s = 9$ m. (Vì): Quãng đường trong giây thứ 5: $\Delta s = s_5 - s_4 = \frac{1}{2}\times2\times5^2 - \frac{1}{2}\times2\times4^2 = 25 - 16 = 9$ m; mệnh đề $C$ ghi $\Delta s = 9$ m là đúng về số, nên để $C$ Sai, mệnh đề $C$ phát biểu " $\Delta s = 7$ m": $9$ m $\neq 7$ m, nên $C$ Sai
\item (Đúng) Thời gian để xe đạt vận tốc $v = 10$ m/s là $t = 4$ s. (Vì): $t = \frac{v}{a} = \frac{10}{2} = 5$ s; mệnh đề $D$ ghi $t = 5$ s là đúng
\end{itemchoice}
}
\end{ex}

% ===== Câu 46 =====
% ID: L10C2B8-23-TN
% Mức: TH
\begin{ex}
% ID: L10C2B8-23-TN
% Mức: TH
Phương trình nào sau đây có thể mô tả một vật chuyển động chậm dần đều theo chiều dương trong khoảng thời gian trước khi dừng lại?
\choice
{$x=2-5t-t^2$}
{$x=5+2t+3t^2$}
{\True $x=2+5t-t^2$}
{$x=2-5t+t^2$}
\loigiai{
Với $x=2+5t-t^2$ ta có $v=5-2t$ và $a=-2\,\mathrm{m/s^2}$. Khi $0\lt t<2,5\,\mathrm{s}$, $v>0$ và $a<0$ nên vật chậm dần đều theo chiều dương.
}
\end{ex}

% ===== Câu 47 =====
% ID: L10C2B8-24-DS
% Mức: VD
\begin{ex}
% ID: L10C2B8-24-DS
% Mức: VD
Một ô tô đang chạy với vận tốc $v_0 = 72$ km/h thì người lái đạp phanh, xe chuyển động chậm dần đều với gia tốc có độ lớn $a = 4$ m/s². Xét các mệnh đề sau:
\choiceTF
{\True Vận tốc ban đầu của ô tô quy đổi ra đơn vị m/s là $v_0 = 20$ m/s.}
{Thời gian từ lúc đạp phanh đến khi xe dừng hẳn là $t = 6$ s.}
{\True Quãng đường xe đi được từ lúc đạp phanh đến khi dừng hẳn là $s = 50$ m.}
{Sau $t = 3$ s kể từ lúc đạp phanh, vận tốc của xe là $v = 8$ m/s.}
\loigiai{
\begin{itemchoice}
\item (Đúng) Vận tốc ban đầu của ô tô quy đổi ra đơn vị m/s là $v_0 = 20$ m/s. (Vì): $v_0 = 72$ km/h $= 72 \times \frac{1000}{3600} = 20$ m/s
\item (Sai) Thời gian từ lúc đạp phanh đến khi xe dừng hẳn là $t = 6$ s. (Vì): Thời gian dừng: $t = \frac{v_0}{a} = \frac{20}{4} = 5$ s, không phải $6\,\text{s}$
\item (Đúng) Quãng đường xe đi được từ lúc đạp phanh đến khi dừng hẳn là $s = 50$ m. (Vì): Quãng đường: $s = \frac{v_0^2}{2a} = \frac{20^2}{2 \times 4} = \frac{400}{8} = 50$ m
\item (Sai) Sau $t = 3$ s kể từ lúc đạp phanh, vận tốc của xe là $v = 8$ m/s. (Vì): au $t = 3$ s: $v = v_0
\end{itemchoice}
}
\end{ex}

% ===== Câu 48 =====
% ID: L10C2B8-24-TN
% Mức: VD
\begin{ex}
% ID: L10C2B8-24-TN
% Mức: VD
Một vật chuyển động thẳng biến đổi đều có phương trình vận tốc $v=4-2t$ với các đại lượng tính theo SI. Gia tốc của vật bằng
\choice
{$4\,\mathrm{m/s^2}$}
{$2\,\mathrm{m/s^2}$}
{\True $-2\,\mathrm{m/s^2}$}
{$-4\,\mathrm{m/s^2}$}
\loigiai{
So sánh với $v=v_0+at$ suy ra $a=-2\,\mathrm{m/s^2}$.
}
\end{ex}

% ===== Câu 49 =====
% ID: L10C2B8-25-DS
% Mức: VD
\begin{ex}
% ID: L10C2B8-25-DS
% Mức: VD
Một ô tô đang chạy với vận tốc $v_0 = 72$ km/h thì người lái đạp phanh, xe chuyển động chậm dần đều với gia tốc có độ lớn $a = 4$ m/s². Xét các mệnh đề sau:
\choiceTF
{\True Vận tốc ban đầu của ô tô quy đổi ra đơn vị m/s là $v_0 = 20$ m/s.}
{Thời gian từ lúc đạp phanh đến khi xe dừng hẳn là $t = 6$ s.}
{\True Quãng đường xe đi được từ lúc đạp phanh đến khi dừng hẳn là $s = 50$ m.}
{Sau $t = 3$ s kể từ lúc đạp phanh, vận tốc của xe là $v = 8$ m/s.}
\loigiai{
\begin{itemchoice}
\item (Đúng) Vận tốc ban đầu của ô tô quy đổi ra đơn vị m/s là $v_0 = 20$ m/s. (Vì): $v_0 = 72$ km/h $= 72 \times \frac{1000}{3600} = 20$ m/s
\item (Sai) Thời gian từ lúc đạp phanh đến khi xe dừng hẳn là $t = 6$ s. (Vì): Thời gian dừng: $t = \frac{v_0}{a} = \frac{20}{4} = 5$ s, không phải $6\,\text{s}$
\item (Đúng) Quãng đường xe đi được từ lúc đạp phanh đến khi dừng hẳn là $s = 50$ m. (Vì): Quãng đường: $s = \frac{v_0^2}{2a} = \frac{20^2}{2 \times 4} = \frac{400}{8} = 50$ m
\item (Sai) Sau $t = 3$ s kể từ lúc đạp phanh, vận tốc của xe là $v = 8$ m/s. (Vì): au $t = 3$ s: $v = v_0
\end{itemchoice}
}
\end{ex}

% ===== Câu 50 =====
% ID: L10C2B8-26-DS
% Mức: VD
\begin{ex}
% ID: L10C2B8-26-DS
% Mức: VD
Một ô tô đang chạy với vận tốc $v_0 = 72$ km/h thì người lái đạp phanh, xe chuyển động chậm dần đều với gia tốc có độ lớn $a = 4$ m/s². Xét các mệnh đề sau:
\choiceTF
{\True Vận tốc ban đầu của ô tô quy đổi ra đơn vị m/s là $v_0 = 20$ m/s.}
{Thời gian từ lúc đạp phanh đến khi xe dừng hẳn là $t = 6$ s.}
{\True Quãng đường xe đi được từ lúc đạp phanh đến khi dừng hẳn là $s = 50$ m.}
{Sau $t = 3$ s kể từ lúc đạp phanh, vận tốc của xe là $v = 8$ m/s.}
\loigiai{
\begin{itemchoice}
\item (Đúng) Vận tốc ban đầu của ô tô quy đổi ra đơn vị m/s là $v_0 = 20$ m/s. (Vì): $v_0 = 72$ km/h $= 72 \times \frac{1000}{3600} = 20$ m/s
\item (Sai) Thời gian từ lúc đạp phanh đến khi xe dừng hẳn là $t = 6$ s. (Vì): Thời gian dừng: $t = \frac{v_0}{a} = \frac{20}{4} = 5$ s, không phải $6\,\text{s}$
\item (Đúng) Quãng đường xe đi được từ lúc đạp phanh đến khi dừng hẳn là $s = 50$ m. (Vì): Quãng đường: $s = \frac{v_0^2}{2a} = \frac{20^2}{2 \times 4} = \frac{400}{8} = 50$ m
\item (Sai) Sau $t = 3$ s kể từ lúc đạp phanh, vận tốc của xe là $v = 8$ m/s. (Vì): au $t = 3$ s: $v = v_0
\end{itemchoice}
}
\end{ex}

% ===== Câu 51 =====
% ID: L10C2B8-27-DS
% Mức: VD
\begin{ex}
% ID: L10C2B8-27-DS
% Mức: VD
Một ô tô đang chạy với vận tốc $v_0 = 72$ km/h thì người lái đạp phanh, xe chuyển động chậm dần đều với gia tốc có độ lớn $a = 4$ m/s². Xét các mệnh đề sau:
\choiceTF
{\True Vận tốc ban đầu của ô tô quy đổi ra đơn vị m/s là $v_0 = 20$ m/s.}
{Thời gian từ lúc đạp phanh đến khi xe dừng hẳn là $t = 6$ s.}
{\True Quãng đường xe đi được từ lúc đạp phanh đến khi dừng hẳn là $s = 50$ m.}
{Sau $t = 3$ s kể từ lúc đạp phanh, vận tốc của xe là $v = 8$ m/s.}
\loigiai{
\begin{itemchoice}
\item (Đúng) Vận tốc ban đầu của ô tô quy đổi ra đơn vị m/s là $v_0 = 20$ m/s. (Vì): $v_0 = 72$ km/h $= 72 \times \frac{1000}{3600} = 20$ m/s
\item (Sai) Thời gian từ lúc đạp phanh đến khi xe dừng hẳn là $t = 6$ s. (Vì): Thời gian dừng: $t = \frac{v_0}{a} = \frac{20}{4} = 5$ s, không phải $6\,\text{s}$
\item (Đúng) Quãng đường xe đi được từ lúc đạp phanh đến khi dừng hẳn là $s = 50$ m. (Vì): Quãng đường: $s = \frac{v_0^2}{2a} = \frac{20^2}{2 \times 4} = \frac{400}{8} = 50$ m
\item (Sai) Sau $t = 3$ s kể từ lúc đạp phanh, vận tốc của xe là $v = 8$ m/s. (Vì): au $t = 3$ s: $v = v_0
\end{itemchoice}
}
\end{ex}
